% !TeX spellcheck = en_GB
%!TEX root = ../main/main.tex
In this section we introduce an extension to the semantics of CEL, namely we introduce a new operator using \textit{Allen interval algebra} \textit{starts}.
We then extend the formal computational model for evaluating CEL formulas and prove its correctness.

We present then the \emph{starts operator} for CEL with the following definition:


$$\begin{aligned} 	\sem{\varphi_1 \starts \varphi_2}(\Stream) 	&= \{C_1 \cup C_2 \mid C_1 \in \sem{\varphi_1}(\Stream)  	\land C_2 \in \sem{\varphi_2}(\Stream) \\ 	&\quad \land \cstart(C_1) = \cstart(C_2) \land \cend(C_1) < \cend(C_2)\} \end{aligned}$$

\begin{example} Suppose we have a system monitoring server health, specifically tracking system errors ($E$), CPU stabilization ($C$), and Memory stabilization ($M$). Consider a stream $\Stream = e_1 e_2 e_3$ where the event types are $\type(e_1) = E$, $\type(e_2) = C$, and $\type(e_3) = M$.

Let $\varphi_1 = E ; C$ be a formula detecting an error followed by CPU stabilization, and $\varphi_2 = E ; M$ be a formula detecting an error followed by Memory stabilization.
\begin{itemize}
	\item The set $\sem{\varphi_1}(\Stream)$ contains a complex event $C_1$ representing the error-to-CPU recovery sequence at positions $1$ and $2$, with $\cinterval(C_1) = [1..2]$. Thus, $\cstart(C_1) = 1$ and $\cend(C_1) = 2$.
	\item The set $\sem{\varphi_2}(\Stream)$ contains a complex event $C_2$ representing the error-to-Memory recovery sequence at positions $1$ and $3$, with $\cinterval(C_2) = [1..3]$. Thus, $\cstart(C_2) = 1$ and $\cend(C_2) = 3$.
\end{itemize}

Since the starts condition $\cstart(C_1) = \cstart(C_2)$ and $\cend(C_1) < \cend(C_2)$ holds (i.e., $1 = 1$ and $2 < 3$), the complex event $C = C_1 \cup C_2$ is in the result of $\sem{\varphi_1 \starts \varphi_2}(\Stream)$. This resulting complex event captures a scenario where both recovery processes were triggered simultaneously by the exact same error, but the CPU stabilized strictly before the Memory did.
\end{example}

We denote by $\celplus{\startssimb}$ the \emph{extension of CEL} with the starts operation. We can prove the following result:

\begin{theorem} \label{theo:mytheorem}
	For every $\celplus{\startssimb}$ formula $\varphi$ there exists a CEA $\cA_\varphi$ such that $\sem{\varphi}(\Stream) = \sem{\cA_\varphi}(\Stream)$ for every stream $\Stream$.
\end{theorem}

\begin{proof}
	We already know that for every (standard) CEL operator we can construct an equivalent CEA. So, we concentrate next in the construction of a CEA for the starts operator.
	Let $\varphi_1$ and $\varphi_2$ be formulas in CEL extended with the starts operator. We show by induction that there exists a CEA $\cAstarts$ such that:
	
	
	$$\sem{\varphi_1 \starts \varphi_2}(\Stream) = \sem{\cAstarts}(\Stream)$$
	
	
	Assume then that there exist CEAs that satisfy the previous property for $\varphi_1$ and $\varphi_2$, called $\cA_{\varphi_1} = (Q_1, \SETP_1, \SETX_1, \Delta_1, q_0, F_1)$ and $\cA_{\varphi_2} = (Q_2, \SETP_2, \SETX_2, \Delta_2, p_0, F_2)$, respectively.
	Then the construction for the starts operator is as follows.
	Let $\cAstarts$ be a CEA where $\cAstarts = (\Qstarts, \Pstarts, \Xstarts, \Dstarts, I_{\starts}, F_{\starts})$:
	
	
	$$\begin{aligned}[t] 		& \Qstarts = Q_2 \uplus (Q_1 \times Q_2 \times \{0, 1\}) \\ 		& \Pstarts = \SETP_2 \cup \{ P_1 \land P_2 \mid P_1 \in \SETP_1 \land P_2 \in \SETP_2 \} \hspace*{155px}\\ 		& \Xstarts  = \SETX_1 \cup \SETX_2 \\ 		& I_{\starts} = (q_0, p_0, 0) \text{ where } q_0 \text{ and } p_0 \text{ are the initial states of } \cA_{\varphi_1} \text{ and } \cA_{\varphi_2}.\\ 		& F_{\starts} = F_2 \text{ the final states of } \cA_{\varphi_2}. 	\end{aligned}$$
	
	
	The transition relation $\Dstarts$ is formally defined as:
	
	
	$$\begin{aligned}[t] 		\Dstarts & = \Delta_2\\ 		& \uplus \{((q, p, b), P_1 \land P_2, L_1 \cup L_2, (q', p', 1)) \mid (q, P_1, L_1, q') \in \Delta_1 \land (p, P_2, L_2, p') \in \Delta_2 \land b \in \{0, 1\}\} \\ 		& \uplus \{((q, p, b), \epsilon, (q', p, b)) \mid (q, \epsilon, q') \in \Delta_1 \land p \in Q_2 \land b \in \{0, 1\}\}  \\ 		& \uplus \{((q, p, b), \epsilon, (q, p', b)) \mid q \in Q_1 \land (p, \epsilon, p') \in \Delta_2 \land b \in \{0, 1\}\}  \\ 		& \uplus \{((q, p, 1), \epsilon, p) \mid q \in F_1 \land p \in Q_2\} 	\end{aligned}$$
	
	Intuitively, given a stream $\Stream$, we evaluate $\varphi_1$ and $\varphi_2$ simultaneously from the beginning. At some point when $\varphi_1$ reaches a final state (ensuring at least one joint event was consumed), we transition out of the product phase and continue evaluating only $\varphi_2$ until it reaches its own final state.
	
	Next, we prove that the construction is correct, namely, $\sem{\varphi_1 \starts \varphi_2}(\Stream) = \sem{\cAstarts}(\Stream)$ for every stream $\Stream$. The proof is by double containment.
	
	\paragraph{$\sem{\varphi_1 \starts \varphi_2}(\Stream) \subseteq \sem{\cAstarts}(\Stream)$}
	
	\paragraph{$\sem{\cAstarts}(\Stream) \subseteq \sem{\varphi_1 \starts \varphi_2}(\Stream)$}
\end{proof}